\documentclass[12pt]{article}
\usepackage[utf8]{inputenc}
\usepackage[T1]{fontenc}
\usepackage[margin=1in]{geometry}
\usepackage{setspace}
\usepackage{xcolor}
\usepackage[colorlinks=true, urlcolor=blue!60!black, linkcolor=blue!60!black]{hyperref}
\setlength{\parindent}{0em}
\setlength{\parskip}{0.6em}
\renewcommand{\baselinestretch}{1.0}
\begin{document}
\begin{center}
{\LARGE \textit{Zivotofsky v.\ Kerry}, 576 U.S.\ 1 (2015)\par}
\vspace{0.5em}
{\large Case Brief --- Two-Page Edition\par}
\end{center}

\section{Memory Jogger}
The President alone recognizes foreign sovereigns, so Congress could not require ``Israel'' on Jerusalem-born citizens' passports (\href{https://supreme.justia.com/cases/federal/us/576/1/}{Justia}).

\section{Detailed Case Facts}
Every President since 1948 has declined to recognize any nation's sovereignty over Jerusalem, so passports list only ``Jerusalem.'' Section 214(d) (2002) required the Secretary, on request, to record ``Israel.'' Menachem Zivotofsky's parents sought that listing and were refused.

\section{Procedural History}
After \textit{Zivotofsky I} (2012) held the case justiciable, the D.C.\ Circuit held \S~214(d) unconstitutional.

\section{Judicial Votes}
Affirmed. Kennedy wrote for Ginsburg, Breyer, Sotomayor, and Kagan; Breyer concurred; Thomas concurred in part and dissented in part; Roberts (with Alito) and Scalia (with Roberts and Alito) dissented.

\section{Holding}
The recognition power is exclusive to the President, and Congress may not compel the Executive to contradict a recognition determination in a passport.

\section{Analysis of Opinions}
\textbf{Majority:} In \textit{Youngstown} Category Three, the Reception Clause, the President's diplomatic role, the need to speak with one voice, precedent, and historical acquiescence establish exclusivity, though the Court rejected \textit{Curtiss-Wright}'s ``sole organ'' dicta. \textbf{Thomas:} the residual Vesting Clause power covers passports but not consular birth reports. \textbf{Roberts:} a passport entry is not recognition, and the ruling allows an ``international heckler's veto.'' \textbf{Scalia:} moving from legal acts to statements is ``a leap worthy of the Mad Hatter.''

\section{Examples: Future Applications}
\textbf{Same-Side:} a law compelling official maps that contradict recognition. \textbf{Opposite-Side:} trade sanctions on a recognized government under the commerce power. \textbf{Fence-Sitter:} customs labels naming an unrecognized republic.

\section{Critique}
The opinion mixes text and functionalism, and its reach over ``official executive statements'' is unclear; \textit{Trump v.\ United States} later cited it for exclusive presidential power.

\section{Key Quotations}
\begin{quote}
``This case is confined solely to the exclusive power of the President to control recognition determinations, including formal statements by the Executive Branch acknowledging the legitimacy of a state or government and its territorial bounds. Congress cannot command the President to contradict an earlier recognition determination in the issuance of passports.''
\end{quote}
\noindent --- Justice Kennedy, for the Court, 576 U.S.\ at 32.

\end{document}
